% ============================================================
%  ch03b  §3.2 同余关系
%  原文对应：ICM-arXiv.tex 1107--1140 行
%  编号：公式 (3.2)；定理 3.2.1；定理 3.2.2
% ============================================================

\subsection{同余关系}\label{S: Congruence}

沿用 \S\ref{geometric JL} 中关于志村簇的记号，并（为简单起见）保留其中的假设。设 $(G,X)$ 是阿贝尔型志村数据，$K$ 是满足 $K_p$ 超特殊的水平。设 $v\mid p$ 是 $E$ 的位置。则 $\Sh_K(G)$ 有定义在 $\mO_{E,(v)}$ 上的典范整模型 $\scrS_K$（\cite{Kisin-Integral-Model}）。

以 $\overline{\scrS}_K$ 记其 mod $p$ 纤维，它是定义在 $v$ 的剩余域 $k_v$ 上的光滑簇。以 $\sigma_v$ 记 $\Ga_{k_v}$ 中的几何 Frobenius。定理 \ref{introT: spectral action} 给出 $\End_{\Loc_{\CG,p}^{\ur}}(\widetilde{V_\mu})$ 在 $H^*_c(\overline\scrS_{K,\overline{k_v}},\bZ_\ell)$ 上的作用，下文将看到它有重要的推论。

同余关系猜想（又称 Blasius--Rogawski 猜想）推广了模曲线的经典 Eichler--Shimura 同余关系 $\Fr_p=T_p+V_p$；它预测：在 $\overline{\scrS}_K\times \overline{\scrS}_K$ 的 Chow 群中，$\overline{\scrS}_K$ 的 Frobenius 自同态满足一个多项式方程，其系数是某些 Hecke 对应的 mod $p$ 约化。定理 \ref{introT: spectral action} 与 \cite[Sect. 6.3]{Xiao-Zhu-Invariants} 一起，在上同调层次上蕴含了这一猜想。

对 ${}^c(G_{\bQ_p})$ 的每个表示 $V$，其特征 $\chi_V$（看作 $\CG|_{d=(p,\sigma_p)}$ 上的 $\hat{G}$-不变函数）通过 Satake 同构 \eqref{Sat isom} 给出元素 $h_V\in H_{G(\bZ_p)}^{\cl}$。

\begin{thm}
下述恒等式
\begin{equation}\label{E: congruence relation}
\sum_{i=0}^{n}(-1)^j h_{\chi_{\wedge ^j V}} \sigma_v^{\dim V-j}=0
\end{equation}
在 $\End (H^*_c(\overline{\scrS}_{K ,\overline{k_v}},\bZ_\ell))$ 中成立，其中 $V=\on{Ind}_{{}^c(G_{E_v})}^{{}^c(G_{\bQ_p})}V_\mu$ 是 $V_\mu$ 的张量诱导\zy{原文 \eqref{E: congruence relation} 中求和指标写作 $i$，而符号与指数处写作 $j$，且未说明上界 $n$ 的取值。此处照录原文，未作改动；疑为 $i$、$j$ 混用之误，已登记于 errata 表。}。
\end{thm}

事实上，由 \cite[Sect. 6.3]{Xiao-Zhu-Invariants}，把 $h_{\chi_{\wedge ^i V}}$ 换成 $S(\chi_{\wedge ^i V})$（其中 $S$ 来自定理 \ref{introT: spectral action} 的 \eqref{introT: spectral action-1}）后等式成立；再由该定理的第 \eqref{introT: spectral action-2} 部分，便可把 $S(\chi_{\wedge ^i V})$ 换回 $h_{\chi_{\wedge ^i V}}$。这条通向 \eqref{E: congruence relation} 的途径，是 V. Lafforgue 对 $\Sht_K(G)$ 的 Eichler--Shimura 关系之处理方式 \cite{Lafforgue-Automorphic-to-Galois} 在志村簇上的类比。传统上还有另一条途径，即直接研究 Hecke 算子的 mod $p$ 约化，Siegel 模簇情形自 \cite{Faltings-Chai-Abelian} 开始；最新进展及相关文献见 \cite{Lee-ES-Relation}。那条途径会在代数对应的层次上给出 \eqref{E: congruence relation}。

今设 $(G,X)=(\Res_{F^+/\bQ}(G_0)_{F^+}, \prod_{\varphi :F^+\to \bR} X_0)$，其中 $(G_0,X_0)$ 是志村数据，$F^+$ 是全实域。如前，取素数 $p$ 使 $K_p$ 超特殊，特别地 $p$ 在 $F^+$ 中非分歧。此外为简单起见，设 $G_{0,\bQ_p}$ 分裂（于是对 $E$ 上位于 $p$ 之上的位置 $v$，有 $E_v=\bQ_p$）。以 $\bF$ 记 $\bF_p$ 的代数闭包。

设 $\{w_i\}_i$ 是 $F^+$ 上位于 $p$ 之上的素数集合，$k_i$ 为 $w_i$ 的剩余域。对每个 $i$，再固定嵌入 $\rho_i: k_i\to\bF$。于是有自然映射
\[
\prod_i(\bZ^{f_i}\rtimes \frakS_{f_i})\to \End_{\Loc_{\CG,p}^{\ur}}(\widetilde{V_\mu}),
\]
其中 $\frakS_{f_i}$ 是 $f_i$ 个文字上的置换群。结合定理 \ref{introT: spectral action}，便得到下述结果（\cite{Xiao-Zhu-Cycle}）。

\begin{thm}
存在 $\prod_i(\bZ^{f_i}\rtimes \frakS_{f_i})$ 在 $H_c^*(\overline\scrS_{K,\overline{\bF}},\bZ_\ell)$ 上的作用，使得 $\sigma_p$ 的作用分解为 $\sigma_p=\prod_i \sigma_{p,i}$，其中 $\sigma_{p,i}=((1,0,\ldots,0), (12\cdots f_i))\in \bZ^{f_i}\rtimes\frakS_{f_i}$。每个 $\sigma_{p,i}^{f_i}$ 都满足一个类似于 \eqref{E: congruence relation} 的多项式方程。
\end{thm}

该定理在某种程度上反映了 Nekovář--Scholl 的 plectic 上同调猜想 \cite{Nekovar-Scholl-Plectic}。
